\section{A local change of variables and cancellation of mixed critical terms}
\label{sec:theta-orbital-identity}

We continue with the same local data as in the preceding section:
$K$ is a finite product of finite extensions of $\mathbb Q_p$,
$R\subset K$ is a compact open subring, and $\psi$ is a continuous
additive character whose annihilator of $R$ is $R$ itself.  Haar
measure is normalized by $\operatorname{vol}(R)=1$.  This includes
the nonmaximal split order already described.  Absolute values and
small neighborhoods in $K$ are understood componentwise.  In
particular, a parameter tends to zero in every field factor.

Fix an integer $m\ge1$, put $s=m+1$, and let $\Omega\subset K^\times$
be compact and open.  Suppose that $\Omega$ is invariant under
multiplication by every $s$th root of unity in $K$.
For $t\in K^\times$, define
\begin{equation}\label{orb-cut}
 \mathcal C_t=\{(-t)^{-m}z^s:z\in\Omega\}.
\end{equation}
Fix an additive lattice $\Lambda$ such that
$\Omega+\Lambda=\Omega$.  A function on $\Omega$ is called
$\Lambda$-invariant when $f(z+\lambda)=f(z)$ for
$z\in\Omega$ and $\lambda\in\Lambda$.
We may always replace $\Lambda$ by a smaller lattice.  None of the
neighborhoods below depends on the values of $f$ or $g$.

\subsection{Uniform choices near one}

We record the elementary local facts used to keep the dependence on
residue characteristic and conductor explicit.

\begin{lemma}\label{orb-near-one}
For each positive integer $s$, the power series
$\beta_s(X)=(1+X)^{1/s}$ converges on a sufficiently small
neighborhood of zero in $K$.  On a smaller neighborhood it satisfies
\[
 \beta_s(X)^s=1+X,\qquad
 \beta_s(X)=1+\frac{X}{s}b_s(X),\qquad b_s(X)\in R,
 \quad b_s(0)=1.
\]
It is the unique $s$th root sufficiently close to one.
Given $\Omega$ and $\Lambda$ as above, there is a neighborhood
$U$ of one in $K^\times$, invariant under inversion, such that
\[
 u\Omega=\Omega,\qquad (u-1)\Omega\subset\Lambda,
 \qquad |u_j|=1
\]
for every $u\in U$ and every field component $j$.
\end{lemma}

\begin{proof}
In a field factor the coefficients of the binomial series have at
most exponential growth.  Indeed, for a fixed rational number $a$,
\[
 \left|\binom a n\right|
 \le\max(1,|a|)^n|n!|^{-1},
\]
and the valuation formula for $n!$ gives an exponential upper bound
for its inverse absolute value.  The series therefore converges on
a sufficiently small ball.  Its formal identity
$\beta_s(X)^s=1+X$ is an identity of convergent series there.
The function $b_s(X)=s(\beta_s(X)-1)/X$, extended by $b_s(0)=1$,
is continuous.  Since $R$ is a neighborhood of one, it takes
values in $R$ after another shrink.  Two nearby roots differ by an
$s$th root of unity.  There are finitely many such roots, and one
can choose the neighborhood to exclude all but one.  Applying these
facts in the finitely many field factors proves the assertion for
$K$.

For the second assertion, compactness of $\Omega$ and continuity
of multiplication give $(u-1)\Omega\subset\Lambda$ for every
$u$ in a sufficiently small neighborhood of one.  Then
$u\Omega\subset\Omega$.  Intersect this neighborhood with its
inverse to obtain equality.  The component absolute values of a
unit near one are all one.
\end{proof}

The neighborhood in this lemma is allowed to depend on $s$, on the
fields, and on $R$.  In particular, it is not asserted to be
independent of the valuations of $s$ or of $2$.  The lemma remains
valid when the residue characteristic divides either integer.

Let $\Sigma\subset K$ be finite.  For $t$ sufficiently small and
$\sigma\in\Sigma$, set
\begin{equation}\label{orb-shift}
 t_\sigma=(t^{-1}+\sigma)^{-1}=\frac{t}{1+\sigma t}.
\end{equation}
The denominator is then a unit, $t_\sigma/t$ tends uniformly to one,
and $|t_\sigma|=|t|$.  Taking the near-one $s$th root of
$(t_\sigma/t)^m$ and using Lemma~\ref{orb-near-one} gives
\begin{equation}\label{orb-cut-shift}
 \mathcal C_{t_\sigma}=\mathcal C_t
 \qquad(\sigma\in\Sigma).
\end{equation}
The same argument shows that $\mathcal C_t+R=\mathcal C_t$
when $t$ is small: if $c=(-t)^{-m}z^s$ and $r\in R$, then
$r/c$ tends uniformly to zero, and
$z\beta_s(r/c)\in\Omega$ is an $s$th root of $(-t)^m(c+r)$.
Compactness of $R$ and of the set of inverses of elements of
$\Omega$ makes this uniform in $c$ and $r$.

\subsection{An exact formula for the sum over root fibers}

Choose any representatives $c_x\in K^\times$ of the additive
cosets $x\in K/R$.  Such representatives exist: the compact open
ring contains a product of sufficiently small open balls, so each
coordinate of a representative can be perturbed away from zero.
The following identity is independent of these choices.

\begin{theorem}\label{orb-change-variable}
For all sufficiently small $t\in K^\times$ and every
$\Lambda$-invariant function $f:\Omega\to\mathbb C$,
\begin{equation}\label{orb-change-formula}
 \sum_{\substack{x\in K/R\cr c_x\in\mathcal C_t}}
       \ \sum_{z^s=(-t)^m c_x} f(z)
 =\frac{|s|}{|t|^m}\int_\Omega |z|^m f(z)\,dz.
\end{equation}
All roots in the inner sum belong to $\Omega$, and both sums are
finite.  The required small neighborhood depends only on
$K,R,m,\Omega,\Lambda$.
\end{theorem}

\begin{proof}
If one root of $z^s=(-t)^mc$ lies in $\Omega$, every root does,
because their ratios are $s$th roots of unity.  Put
\[
 H(c)=\sum_{\substack{z\in\Omega\cr z^s=(-t)^mc}}f(z).
\]
This function has compact support $\mathcal C_t$.  For
$c\in\mathcal C_t$ and $r\in R$, multiplication by
$\omega=\beta_s(r/c)$ bijects the roots above $c$ and $c+r$.
For small $t$, uniformly over these choices,
$z(\omega-1)\in\Lambda$.  Thus it preserves $f$, and
$H(c+r)=H(c)$.  The invariance of $\mathcal C_t$ under $R$
proves the same equality outside $\mathcal C_t$.  Consequently
\[
 \sum_{x\in K/R}H(c_x)=\int_K H(c)\,dc,
\]
since each additive coset has volume one and only finitely many
cosets meet the compact support.

The map $z\mapsto c=z^s/(-t)^m$ has nonzero derivative
$s z^m/(-t)^m$ at every point of $\Omega$.  On sufficiently small
product balls it is an analytic bijection onto a product of balls,
and multiplies Haar measure by $|s||z|^m/|t|^m$.  To justify this
without a smoothness assumption on $f$, divide the local map by
its derivative: on a smaller ball it is a linear map followed by
the identity plus a strict contraction, exactly as in the proof
of the local Morse lemma.  It therefore maps each sufficiently
small coset partition bijectively to the corresponding partition.
The Jacobian factor is constant on that ball.  A finite disjoint
partition of $\Omega$ by such balls proves the change-of-variables
formula, with every root counted once.  This gives
\eqref{orb-change-formula}.
\end{proof}

The source writes the integral on the right as a finite lattice
sum and additionally assumes that $|z|$ is constant on the
$\Lambda$-cosets in $\Omega$.  Under that assumption its sum is
exactly the Haar integral in \eqref{orb-change-formula}.  The
Haar-integral version itself needs no such extra assumption.

\subsection{Vanishing of the mixed critical terms}

\begin{theorem}\label{orb-cross-vanish}
Fix $\zeta\in K^\times$ with $\zeta^s=1$ and $\zeta\ne1$,
and fix a finite set $\Sigma\subset K$.  There are neighborhoods
$E$ of one and $V$ of zero such that the following holds.
Let $t\in V\cap K^\times$, let
$t_1=t_\sigma$, $t_2=t_{\sigma'}$ with
$\sigma,\sigma'\in\Sigma$, and let $\kappa\in E$.
Suppose that $g:\Omega\to\mathbb C$ is $\Lambda$-invariant
and satisfies
\begin{equation}\label{orb-cross-support}
 g(z)\ne0\quad\Longrightarrow\quad
 z^{-1}-(\zeta z)^{-1}\in R.
\end{equation}
Then, for every choice of the representatives $c_x$,
\begin{equation}\label{orb-cross-sum}
 \sum_{\substack{x\in K/R\cr c_x\in\mathcal C_t}}
 \ \sum_{z^s=(-t_1)^m c_x}
 g(z)\psi\left(sz/t_1-s\zeta\kappa z/t_2\right)=0.
\end{equation}
In particular this holds with the source's further condition
$\kappa^s=(t_2/t_1)^m$; that condition is unnecessary for this
vanishing assertion alone.
\end{theorem}

\begin{proof}
Write
\[
 a=\frac{s}{t_1}-\frac{s\zeta\kappa}{t_2},\qquad
 \rho=\kappa\frac{t_1}{t_2},\qquad
 T=(-t_1)^m.
\]
We may make $\rho$ as close to one as needed by shrinking $E$
and $V$.  The truncation sets for $t_1$ and $t$ agree by
\eqref{orb-cut-shift}.  Define
\[
 H(c)=\sum_{\substack{z\in\Omega\cr z^s=Tc}}
             g(z)\psi(az).
\]
The important point is that $H$ is still $R$-periodic, despite
the oscillating factor.  For $r\in R$, multiplication by
$\omega=\beta_s(r/c)$ matches the roots above $c$ and $c+r$,
preserves $\Omega$, and preserves $g$.  Moreover
\[
 az(\omega-1)=\frac{aT}{s z^m}\,r\,b_s(r/c),\qquad
 \frac{aT}{s z^m}
    =(-1)^m\frac{t_1^{m-1}}{z^m}(1-\zeta\rho).
\]
The second expression belongs to $R$ on the support of $g$,
uniformly for sufficiently small parameters.  If $m\ge2$, it
tends uniformly to zero, since $\Omega$ is compact in $K^\times$.
If $m=1$, then $s=2$ and $\zeta^{-1}=\zeta$, so
\[
 -\frac{1-\zeta\rho}{z}
 =-\frac{1-\zeta}{z}+\frac{\zeta(\rho-1)}{z}\in R.
\]
The first term belongs to $R$ by \eqref{orb-cross-support};
the second does so by the uniform smallness of $\rho-1$.
Since $r,b_s(r/c)\in R$ and $\psi$ is trivial on $R$, the
phase is unchanged.  Terms where $g(z)=0$ cause no difficulty,
because the matched value of $g$ is also zero.  Therefore
$H(c+r)=H(c)$, including outside its $R$-stable support.

As in the preceding theorem, the sum in \eqref{orb-cross-sum}
is consequently
\begin{equation}\label{orb-cross-integral}
 \frac{|s|}{|t_1|^m}
       \int_\Omega |z|^m g(z)\psi(az)\,dz.
\end{equation}
Choose a fixed smaller lattice $\Lambda_0\subset\Lambda$
such that $\Omega$ is a finite union of its cosets and $|z|$
is constant on each such coset.  This is possible because the
absolute value is locally constant on $K^\times$ and $\Omega$
is compact.  The amplitude $|z|^m g(z)$ is constant on each coset,
independently of the values of $g$.

There is at least one field component $j$ with $\zeta_j\ne1$.
After shrinking $E$ and $V$, the number $1-\zeta_j\rho_j$
has the same nonzero absolute value as $1-\zeta_j$.
Thus $|a_j|=|s(1-\zeta_j\rho_j)/t_{1,j}|$ tends to infinity.
The annihilator $\Lambda_0^\perp$ is compact, so
$a\notin\Lambda_0^\perp$ once $t$ is sufficiently small.
Character orthogonality on each coset of $\Lambda_0$ now makes
the integral in \eqref{orb-cross-integral} zero.  This proves
\eqref{orb-cross-sum}.
\end{proof}

The proof isolates the role of the support condition: it is needed
for invariance under a change of representative when $m=1$.
For $m\ge2$ the same conclusion holds without that condition.
All choices remain local choices made after $K,R,\psi,m,\Omega,
\Lambda,\Sigma,\zeta$ are fixed.  No constant uniform over
all residue characteristics or over all block sizes has been
asserted.

\begin{remark}[The support condition cannot simply be omitted when $m=1$]
Take $K=\mathbb Q_3$, $R=\mathbb Z_3$, and a character with
annihilator $R$.  Let $m=1$, $s=2$, $\zeta=-1$,
$\Sigma=\{0\}$, $\kappa=1$, and set
\[
 \Omega=(3+9\mathbb Z_3)\cup(-3+9\mathbb Z_3),\qquad
 \Lambda=9\mathbb Z_3,\qquad g=1_{3+9\mathbb Z_3}.
\]
All the assumptions of Theorem~\ref{orb-cross-vanish}, other than
\eqref{orb-cross-support}, hold.  Put $t=3^N$ with $N\ge4$.
One may choose $c=-9/t$ as the representative of its coset in
$K/R$.  The roots of $z^2=-tc$ are $3$ and $-3$, and the
contribution of this coset is $\psi(12/t)$.
Now replace only this representative by $c+1$.  Its root in
$3+9\mathbb Z_3$ is
\[
 z'=3\sqrt{1-3^{N-2}}.
\]
The near-one square root exists, and its binomial series gives
\[
 \frac{4(z'-3)}{t}
       =-\frac23+O(3^{N-3})
       \equiv-\frac23\pmod{\mathbb Z_3}.
\]
Indeed the first-order term is $-2/3$ and, for $j\ge2$, the
$j$th term has valuation at least $N-3$; the binomial
coefficients of $(1+X)^{1/2}$ lie in $\mathbb Z_3$.
Thus this coset's contribution is multiplied by
$\psi(-2/3)\ne1$, while all other contributions are unchanged.
The two total sums differ, so they cannot both vanish.
Such a failure occurs for arbitrarily small $t$.
The omitted hypothesis fails at $z=3$, since
$z^{-1}-(-z)^{-1}=2/3\notin R$.
This is a counterexample to removing a hypothesis, not to the
statement in the source.
\end{remark}

\subsection{The exact truncated orbital identity and its remaining algebraic inputs}

A block configuration $h$ consists of a finite list of block sizes
and one scalar in $K^\times$ for each block.  Write $h[s,z]$ for
the configuration obtained by adjoining the block $zJ_s$, where
$J_s$ is the reversing permutation matrix of size $s$.  Zero-sized
blocks contribute empty matrices and can be omitted.  The matrix
$Z(h)$ is the block diagonal matrix of the configuration.

Let $\mathcal K(h,h')$ denote the local marked-cokernel kernel of
the preceding section.  For convenience define
\begin{align}
 N_s(h,h';z,z')
   &=|zz'|^{s(s-1)/2}
                 \mathcal K(h[s,z],h'[s,z']),\\
 W_h(t,c)&=h[m,t][1,c].
\end{align}
The truncated Gram sum is
\begin{equation}\label{orb-cut-gram}
 U_t(h,h';t_1,t_2)=|t_1t_2|^{m(m+1)/2}
       \sum_{\substack{x\in K/R\cr c_x\in\mathcal C_t}}
          \mathcal K(W_h(t_1,c_x),W_{h'}(t_2,c_x)).
\end{equation}

The next theorem isolates the exact algebraic inputs still needed
for this particular kernel.  The hypotheses are stated in
ordinary mathematical terms rather than concealed in a general
existence assertion.

\begin{theorem}\label{orb-cut-identity}
Fix a finite collection $\mathcal F$ of block configurations of
the same total size, a compact open $\Omega$ as above, and a
finite $\Sigma\subset K$.  Suppose that the following properties
hold for the kernel $\mathcal K$.
\begin{enumerate}
\item The kernel is Hermitian.
\item For each $h\in\mathcal F$ there is an additive lattice
$L_h$ such that, whenever $z,z'\in K^\times$, $z\in\Omega$,
and $z'-z\in L_h$, one has
\[
 \mathcal K(h[s,z'],w)=\mathcal K(h[s,z],w)
 \quad\text{for every comparison configuration }w.
\]
\item For $h,h'\in\mathcal F$ and $z,z'\in K^\times$,
\begin{equation}\label{orb-inverse-support}
 \mathcal K(h[s,z],h'[s,z'])\ne0
 \quad\Longrightarrow\quad z^{-1}-z'^{-1}\in R.
\end{equation}
\item With $\gamma_m$ the Gaussian phase from the preceding
section, the following germ formula holds.  For every compact
$Q\subset K^\times$ there is a sufficiently small neighborhood
of zero such that, for $h\in\mathcal F$, $t$ in that neighborhood,
$ct^m\in Q$, and every comparison configuration $w$,
\begin{align}\label{orb-germ-input}
 \mathcal K(W_h(t,c),w)
  ={}&|t|^{-m^2/2}|s|^{-1/2}
       \sum_{z^s=(-t)^m c}\gamma_m(zt)|z|^{m^2/2}
             \psi(sz/t)\mathcal K(h[s,z],w).
\end{align}
The neighborhood is uniform in $c$ and $w$ as stated.
\end{enumerate}
Then for all sufficiently small $t\in K^\times$, all
$\sigma,\sigma'\in\Sigma$, and all $h,h'\in\mathcal F$,
with $t_1=t_\sigma$ and $t_2=t_{\sigma'}$, one has
\begin{equation}\label{orb-exact-cut}
 U_t(h,h';t_1,t_2)
   =\int_\Omega N_s(h,h';z,z)
                         \psi(z(\sigma-\sigma'))\,dz.
\end{equation}
The small neighborhood is independent of the representatives $c_x$.
\end{theorem}

\begin{proof}
All choices below involve finitely many configurations, shifts,
and $s$th roots of unity, so the required neighborhoods and lattices
can be intersected.  Choose the unique root near one satisfying
\[
 \kappa^s=(t_2/t_1)^m.
\]
For each root $z^s=(-t_1)^mc$, the roots for the second germ
formula are exactly $\zeta\kappa z$, with $\zeta^s=1$.
The truncation sets for $t,t_1,t_2$ agree.  In particular all these roots
are in $\Omega$.

Take
$Q=\{(-1)^m z^s:z\in\Omega\}$ in the germ hypothesis.
Apply \eqref{orb-germ-input} first to the first argument of the
kernel in \eqref{orb-cut-gram}, and then to its second argument
by Hermitian symmetry.  These are finite sums.  Since
$|t_1|=|t_2|=|t|$ and $|\zeta|=|\kappa|=1$, the result is
\begin{align}\label{orb-expanded-cut}
 U_t(h,h';t_1,t_2)
 ={}&\frac{|t|^m}{|s|}
 \sum_{\zeta^s=1}
 \sum_{\substack{x\in K/R\cr c_x\in\mathcal C_t}}
 \ \sum_{z^s=(-t_1)^m c_x}
 |z|^{-m}\notag\\
 &\quad\cdot\gamma_m(zt_1)
       \overline{\gamma_m(\zeta\kappa zt_2)}
       N_s(h,h';z,\zeta\kappa z)\notag\\
 &\quad\cdot\psi(sz/t_1-s\zeta\kappa z/t_2).
\end{align}
For completeness, the power of $|t|$ comes from
$m(m+1)-m^2=m$.  The power of $|z|$ comes from
$m^2-s(s-1)=-m$.  These equalities account for the whole
normalization in this calculation.

The factor $\kappa t_2/t_1$ is a square when $t$ is sufficiently
small.  Square invariance of the Gaussian phase gives
\[
 \gamma_m(\zeta\kappa zt_2)=\gamma_m(\zeta zt_1).
\]
The local constancy hypothesis, applied in the second argument
using Hermitian symmetry, gives
\[
 N_s(h,h';z,\zeta\kappa z)=N_s(h,h';z,\zeta z).
\]
Here $\zeta z\in\Omega$, and
$(\kappa-1)\zeta z$ is uniformly small.  The positive factor in
the definition of $N_s$ is unchanged because $|\kappa|=1$.

Choose a fixed lattice $\Lambda_0$ small enough that $\Omega$
is a union of its cosets, that $|z|$ is constant on those cosets,
and that $(z+\lambda)/z$ is a square for
$z\in\Omega$, $\lambda\in\Lambda_0$.  Also require
$\Lambda_0\subset L_h$ and $\zeta\Lambda_0\subset L_h$
for every $h\in\mathcal F$ and $\zeta^s=1$.
These choices are possible by compactness and
Lemma~\ref{orb-near-one}.  The amplitude
\[
 g_{\zeta,t_1}(z)=|z|^{-m}\gamma_m(zt_1)
           \overline{\gamma_m(\zeta zt_1)}
                       N_s(h,h';z,\zeta z)
\]
is then $\Lambda_0$-invariant, with that lattice independent of
$t_1$.  Square invariance handles the two Gaussian factors, and
the local constancy hypothesis handles both kernel arguments.
If this amplitude is nonzero, \eqref{orb-inverse-support}
gives exactly \eqref{orb-cross-support}.
Theorem~\ref{orb-cross-vanish} therefore removes every term in
\eqref{orb-expanded-cut} with $\zeta\ne1$.

For $\zeta=1$ the two Gaussian factors multiply to one.
The remaining phase satisfies
\begin{equation}\label{orb-diagonal-phase}
 \psi(sz/t_1-s\kappa z/t_2)
               =\psi(z(\sigma-\sigma'))\qquad(z\in\Omega).
\end{equation}
Here is a direct verification, including the necessary uniformity.
As a convergent power series in the small parameter $t$,
\[
 \kappa=1+\frac{m}{s}(\sigma-\sigma')t+t^2q(t),
\]
where $q$ is analytic near zero.  This follows by substituting
$t_2/t_1=(1+\sigma t)/(1+\sigma't)$ in its defining
near-one root.  Consequently
\[
 s/t_1-s\kappa/t_2
     =\sigma-\sigma'+tq_1(t)
\]
for another analytic function $q_1$.  Its values are bounded on
a smaller neighborhood.  Multiplication by $z\in\Omega$ makes
the remainder lie in $R$ for all sufficiently small $t$.
There are finitely many pairs of shifts, so the choice is uniform.
Triviality of $\psi$ on $R$ proves \eqref{orb-diagonal-phase}.

After shrinking $\Lambda_0$ again to annihilate all differences
$\sigma-\sigma'$, the remaining amplitude
\[
 |z|^{-m}N_s(h,h';z,z)\psi(z(\sigma-\sigma'))
\]
is $\Lambda_0$-invariant.  Theorem~\ref{orb-change-variable}
turns its root sum into
\[
 \frac{|s|}{|t_1|^m}
       \int_\Omega N_s(h,h';z,z)
                      \psi(z(\sigma-\sigma'))\,dz.
\]
This cancels the factor $|t|^m/|s|$ in
\eqref{orb-expanded-cut}, proving \eqref{orb-exact-cut}.
\end{proof}

The theorem proves the full analytic passage from the orbital germ
formula to the truncated orbital identity, including the mixed critical terms.
It does not establish the germ formula or the inverse-parameter
support for an arbitrary Hermitian kernel.  For the actual marked
cokernel kernel those properties are additional algebraic facts;
their exact content is examined next.

\subsection{Finiteness of the actual sums}

We make the kernel definition precise enough for the algebraic
support argument.  For a configuration $w$ of total size $n$, let
$L_w(K)$ be the block lower unipotent group, with identity matrices
on its diagonal blocks and arbitrary blocks strictly below them.
Let $U_n(K)$ be the upper unipotent group.  For either group,
placing $R$ in place of $K$ means that all its entries lie in $R$.
The source's lower group is exactly $L_w(K)$: conjugating a lower
triangular diagonal block by its reversing permutation matrix
makes it upper triangular, so a diagonal block that remains lower
unipotent must be the identity.

For
\[
 q=(x,y)\in L_w(R)\backslash L_w(K)
                    \times U_n(K)/U_n(R),
\]
choose representatives and put
\[
 Z_q=xZ(w)y,\qquad
 a_q=\psi\left(\sum_{i=0}^{n-2}x_{i+1,i}
                       +\sum_{i=0}^{n-2}y_{i,i+1}\right).
\]
Let $\mathcal I_w$ be the set of those cosets for which
$Z_q\in M_n(R)$.  The marked module attached to $Z_q$ is
$C_{Z_q}=R^n/Z_qR^n$, with marked vector $[e_{n-1}]$ and
functional $[v]\mapsto (Z_q^{-1}v)_{n-1}\pmod R$.
For $n=0$, take the zero module with zero markings.  Then the actual kernel is
\begin{equation}\label{orb-actual-kernel}
 \mathcal K(w,w')=
 \sum_{q\in\mathcal I_w}\sum_{q'\in\mathcal I_{w'}}
 a_q\overline{a_{q'}}\,
       \#\operatorname{Iso}(C_{Z_q},C_{Z_{q'}}).
\end{equation}
Isomorphisms here preserve both markings.

\begin{lemma}\label{orb-finite-orbits}
The set $\mathcal I_w$ is finite, and
\eqref{orb-actual-kernel} is independent of every representative
choice.  In particular it is the finite Gram kernel described
earlier.  For $w=h[s,z]$, the union of the integral orbit cosets
is still finite when $z$ ranges over a fixed compact subset of
$K^\times$.
\end{lemma}

\begin{proof}
First consider all pairs $(x,y)$ for which $xZ(w)y$ has entries
in $R$.  Write $xZ(w)y=LU$ in block form, with $L=x$ and
$U=Z(w)y$.  The determinant of the $i$th diagonal block of $U$
is the fixed nonzero number $\det(t_iJ_{s_i})$.
Block Gaussian elimination recovers $L$ and $U$ from the product
$Z$ by
\[
 U_{ij}=Z_{ij}-\sum_{a<i}L_{ia}U_{aj}\quad(j\ge i),\qquad
 L_{ji}=\left(Z_{ji}-\sum_{a<i}L_{ja}U_{ai}\right)U_{ii}^{-1}
 \quad(j>i).
\]
At each step $U_{ii}^{-1}$ is its adjugate divided by its fixed
nonzero determinant.  Thus the entries of $L,U$, and
$y=Z(w)^{-1}U$ are bounded when the entries of $Z$ lie in the
compact set $R$.  All conditions on $x,y$ are closed.  The set
of such pairs is therefore a closed bounded subset of a
finite-dimensional vector space over the local fields, hence
compact.  Its image in the displayed coset space is finite because
the integral unipotent subgroups are open and that coset space
is discrete.  The case $n=0$ has just one pair and the same
conclusion.
The same proof is uniform for $z$ in a compact subset of
$K^\times$: the diagonal-block determinants and their inverses
then range over compact sets.  Apply the recovery formulas on
the compact set of pairs $(Z,z)$ instead.  The group $L_w$
depends only on the block sizes, so the quotient space is fixed.

Changing representatives replaces $Z_q$ by $aZ_qb$, with
$a\in L_w(R)$ and $b\in U_n(R)$.  Both matrices are invertible
over $R$.  The map $[v]\mapsto[av]$ identifies the cokernels.
It preserves the marked vector because $ae_{n-1}=e_{n-1}$,
and preserves the functional because
$e_{n-1}^tb^{-1}=e_{n-1}^t$.
The sums of the adjacent entries are additive on unipotent
groups.  Their changes lie in $R$, so the phase $a_q$ is
unchanged.  This proves representative independence and the
claimed identification.
\end{proof}

\begin{proposition}\label{orb-local-constancy-proved}
Fix $h$, an integer $s\ge1$, and a compact set
$C\subset K^\times$.  There is an additive lattice $L$ such
that for $z\in C$ and $z'\in K^\times$ with $z'-z\in L$,
\[
 \mathcal K(h[s,z'],w)=\mathcal K(h[s,z],w)
 \quad\text{for every configuration }w.
\]
The lattice is independent of the comparison configuration.
\end{proposition}

\begin{proof}
Enlarge $C$ to a compact open subset of $K^\times$ and choose
a preliminary lattice so that its small translates remain
inside this larger set.  By Lemma~\ref{orb-finite-orbits}, a
single finite list of orbit cosets contains all the possible
integral terms.  Choose fixed representatives $(x_q,y_q)$
for them.  Write
$Z_q(z)=x_qZ(h[s,z])y_q$.
These matrices are invertible, and their inverses range over
a compact set as $q,z$ vary.  Choose an integer $b\ge0$
such that every entry of every $Z_q(z)^{-1}$ lies in
$p^{-b}R$.

Shrink the lattice $L$ so that, for every listed coset,
\[
 \Delta_q=Z_q(z')-Z_q(z)\in p^{2b+1}M_n(R)
 \quad\text{whenever }z'-z\in L.
\]
This is possible because the difference is a fixed linear
function of $z'-z$ for each of finitely many $q$.
It implies that $Z_q(z')$ is integral exactly when $Z_q(z)$
is integral.  Put $M=Z_q(z)^{-1}\Delta_q$.
Then $M\in p^{b+1}M_n(R)$, and the geometric series for
$(I+M)^{-1}$ converges inside $M_n(R)$, since $R$ is compact
and hence complete.  Thus
\[
 Z_q(z')R^n=Z_q(z)R^n,
 \qquad Z_q(z')^{-1}-Z_q(z)^{-1}\in pM_n(R).
\]
The first identity identifies the cokernels by the identity
map of $R^n$; the second preserves every inverse-coordinate
functional modulo $R$.  The marked vector is unchanged,
and the phase depends only on $x_q,y_q$, not on $z$.
Consequently every marked module and weight in the finite
sum is unchanged.  Counting isomorphisms to any comparison
marked module gives the asserted equality, uniformly in $w$.
\end{proof}

\subsection{The finite-module character argument}

We use the following elementary duality fact.  Its proof also
explains why a nonmaximal order causes no difficulty.

\begin{lemma}\label{orb-dual-span}
Let $C$ be a finite $R$-module, and let
$\lambda_1,\ldots,\lambda_l,\nu:C\to K/R$ be $R$-linear.
Then $\nu$ is an $R$-linear combination of the $\lambda_i$
if and only if it vanishes on their common kernel.
\end{lemma}

\begin{proof}
Only the reverse implication needs proof.  The character
$c\mapsto\psi(\nu(c))$ then factors through the image of
$\lambda=(\lambda_1,\ldots,\lambda_l)$.
That image is a finite additive subgroup of $(K/R)^l$, contained
in $(p^{-N}R/R)^l$ for a sufficiently large integer $N$.
Every character of a subgroup of a finite abelian group extends
to the group: adjoin one element at a time, choosing a root of
the prescribed value of its least positive multiple lying in
the subgroup.  Self-duality of $R$ gives a perfect pairing
\[
 R/p^NR\ \times\ p^{-N}R/R\longrightarrow\mathbb C^\times,
 \qquad(a,z)\longmapsto\psi(az).
\]
Its kernels are exactly $p^NR$ and $R$; the two finite groups
have the same order by Haar-volume scaling.  Thus the extended
character has the form
$\psi(\sum_i a_i\lambda_i(c))$ with $a_i\in R$.
Apply this equality to $rc$ for every $r\in R$.  Self-duality
then implies
$\nu(c)-\sum_i a_i\lambda_i(c)=0$ in $K/R$ for every $c$.
\end{proof}

For an integral invertible $n$ by $n$ matrix $Z$, put
\[
 \lambda_j^Z([v])=(Z^{-1}v)_j\pmod R
 \qquad(0\le j<n).
\]
These functionals are well-defined on $C_Z$.  If
$\phi:C_Z\to C_{Z'}$ is an isomorphism, let $T_j(\phi)$
mean that for every $i\ge j$ there are $a_{il}\in R$ such that
\begin{equation}\label{orb-triangular}
 \lambda_i^{Z'}\phi=\lambda_i^Z+
                                  \sum_{l>i}a_{il}\lambda_l^Z.
\end{equation}
The marked-module isomorphisms already satisfy $T_{n-1}$.
By Lemma~\ref{orb-dual-span}, assuming $T_{k+1}$, the additional
condition $T_k$ is equivalent to vanishing of
$\lambda_k^{Z'}\phi-\lambda_k^Z$ on the common kernel of the
functionals $\lambda_i^Z$ with $i>k$.

\begin{proposition}\label{orb-triangular-average}
Let $w=h[s,t]$ and $w'=h'[s,t']$, where the two previous
configurations have the same total size $r$, and let
$n=r+s$.  If $r\le k\le n-2$, then the sum in
\eqref{orb-actual-kernel} is unchanged when, among the
isomorphisms counted in each summand, the condition $T_{k+1}$
is strengthened to $T_k$.
\end{proposition}

\begin{proof}
It suffices to show that the weighted contribution of the
isomorphisms satisfying $T_{k+1}$ but not $T_k$ is zero.
We partition the upper-unipotent variables into smaller cosets.
Put $a=k+1$ and $l=n-a$.  Let $P_k\subset U_n(K)$ consist of
matrices whose two diagonal blocks of sizes $a,l$ have entries
in $R$; the upper-right rectangle is unrestricted.  This is a
subgroup containing $U_n(R)$.  Every coset in $P_k/U_n(R)$ has
a representative
\[
 p_B=\begin{pmatrix}I_a&B\\0&I_l\end{pmatrix},
 \qquad B\in M_{a,l}(K)/M_{a,l}(R).
\]
Indeed one can remove both integral diagonal blocks by right
multiplication.  Two displayed representatives define the same
coset exactly when their difference $B-B'$ is integral.

Fix the lower-unipotent cosets and representatives $y_0,y'_0$
of the upper cosets modulo $P_k$.  Write
$Z_0=xZ(w)y_0$ and $Z'_0=x'Z(w')y'_0$.
Only finitely many $B,B'$ make $Z_0p_B,Z'_0p_{B'}$ integral,
by Lemma~\ref{orb-finite-orbits}.  Empty collections contribute
zero and may be discarded.  The first $a$ columns of $Z_0$
and $Z'_0$ are then integral, since multiplication by $p_B$
does not change them.  Define
\[
 S_0=Z_0(R^a\oplus0)\subset R^n,\quad
 P=R^n/S_0,\quad
 \Lambda:P\to K^l,\quad
 \Lambda([v])=\operatorname{pr}_{>k}(Z_0^{-1}v),
\]
and define $S'_0,P',\Lambda'$ in the same way.
The maps are well-defined.  For
$L_B=Z_0p_BR^n\subset R^n$, the restriction
\begin{equation}\label{orb-tail-section}
 \Lambda:L_B/S_0\longrightarrow R^l
\end{equation}
is an isomorphism.  It is surjective by the last $l$ columns
of $Z_0p_B$; its kernel is zero because a vector with last
$l$ inverse coordinates zero belongs to the first $a$ columns.
The analogous statement holds for $L'_{B'}$.

The marked vector $e=e_{n-1}$ satisfies
$\Lambda([e])=\Lambda'([e])=0$.
To check this, use
$Z_0^{-1}e=y_0^{-1}Z(w)^{-1}x^{-1}e
=t^{-1}y_0^{-1}e_r$.
An upper triangular matrix times $e_r$ has zero coordinates
greater than $r$, hence greater than $k$.  The same calculation
applies to $Z'_0$.

We now describe the isomorphisms in this finite sum by data
independent of $B,B'$.  The higher inverse-coordinate maps of
$Z_0p_B$ are just $\Lambda$ modulo $R^l$.
An isomorphism
$\phi:R^n/L_B\to R^n/L'_{B'}$ satisfying $T_{k+1}$ therefore
satisfies
\[
 \Lambda'(v')\equiv D\Lambda(v)\pmod{R^l}
 \quad\text{whenever }\phi([v])=[v'],
\]
for some upper unipotent $D\in U_l(R)$.
Choose once and for all representatives for the equivalence
relation
\[
 D_1\sim D_2\quad\Longleftrightarrow\quad
 (D_1-D_2)\Lambda(P)\subset R^l.
\]
The isomorphism $\phi$ determines exactly one such class.
Replacing $D$ by its chosen representative preserves the
congruence above.

There is a unique lift
\begin{equation}\label{orb-partial-lift}
 \Psi:P\longrightarrow P',\qquad
 \Lambda'\Psi=D\Lambda,
\end{equation}
that induces $\phi$.  Indeed, start with any representative
$v'$ of $\phi([v])$.  Its error
$\Lambda'(v')-D\Lambda(v)$ lies in $R^l$ and can be removed
by subtracting the unique element of $L'_{B'}/S'_0$ with
that image under \eqref{orb-tail-section}.  This defines the
lift independently of all choices and proves linearity.
It is injective because an element in its kernel lies in
$L_B/S_0$ and has zero $\Lambda$-image.  It is surjective by
the same adjustment in the source, using $D^{-1}R^l=R^l$.
The lift carries the marked vector to its counterpart, since
both have zero higher coordinates.

Conversely, a marked isomorphism $\Psi:P\to P'$ satisfying
\eqref{orb-partial-lift}, with $D$ a chosen representative,
descends to exactly one isomorphism $\phi$ whenever
\begin{equation}\label{orb-related-lattices}
 \Psi(L_B/S_0)=L'_{B'}/S'_0.
\end{equation}
Its last inverse-coordinate functional is preserved because
the last row of $D$ is the last coordinate vector; hence it
is an isomorphism of the marked modules.  It satisfies
$T_{k+1}$ by construction.

The failure of $T_k$ is also intrinsic to $\Psi$.  It is
equivalent to the existence of $v,v'\in R^n$ such that
\begin{equation}\label{orb-bad-pair}
 \Psi([v])=[v'],\qquad \Lambda(v)=\Lambda'(v')=0,
 \qquad (Z_0'^{-1}v')_k-(Z_0^{-1}v)_k\notin R.
\end{equation}
To see this, apply Lemma~\ref{orb-dual-span} to the common
kernel of the higher coordinate functionals.  A representative
whose higher coordinates lie in $R^l$ can be changed, using
\eqref{orb-tail-section}, to one whose higher coordinates
are exactly zero.  The corresponding change in the target
is made using $\Psi$.  These changes do not alter any
inverse-coordinate difference modulo $R$.  On exact zero
higher coordinates, multiplication by $p_B^{-1}$ does not
alter the $k$th coordinate, proving \eqref{orb-bad-pair}.
The converse follows by applying the same criterion to that
pair of representatives.

Group all bad isomorphisms by their lift $\Psi$.  Each resulting
set of pairs $(B,B')$ is finite.  Fix such a $\Psi$, choose
$v,v'$ as in \eqref{orb-bad-pair}, and write the last
nonintegral difference there as $d$.  By self-duality there
is $r_0\in R$ with $\psi(r_0d)\ne1$.
Put $v_0=r_0v$, $v'_0=r_0v'$, and
\[
 \xi=Z_0^{-1}v_0,\quad \xi'=Z_0'^{-1}v'_0,
 \quad\beta=(1,0,\ldots,0)\in R^{1\times l},
 \quad\beta'=\beta D^{-1}.
\]
Extend $\beta,\beta'$ by zeros in the first $a$ coordinates.
Both $\xi$ and $\xi'$ vanish in the last $l$ coordinates.
The matrices
\[
 \sigma=I+\xi\beta,\qquad
 \sigma'=I+\xi'\beta'
\]
are therefore upper unipotent shears supported in the
upper-right rectangle.  Their inverses have the corresponding
minus signs.

Replacing $p_B,p_{B'}$ by $\sigma p_B,\sigma'p_{B'}$
preserves integrality: their products with $Z_0,Z'_0$ change
by $v_0\beta,v'_0\beta'$, whose entries lie in $R$.
It also preserves \eqref{orb-related-lattices}.
Here are the details of that assertion.  For $x\in R^n$,
put $b(x)=\beta\Lambda(x)$.  Then
\[
 x\in Z_0\sigma p_BR^n
 \quad\Longleftrightarrow\quad
 b(x)\in R\text{ and }x-b(x)v_0\in L_B.
\]
This follows directly by applying the inverse shear; on
$L_B$ the higher inverse coordinates are integral, and
$\beta\xi=0$.
For representatives $x,y$ with $\Psi([x])=[y]$, the analogous
scalar in the target is
$\beta'\Lambda'(y)=\beta'D\Lambda(x)=b(x)$.
Moreover $\Psi([v_0])=[v'_0]$.  The equivalence of the two
membership conditions now follows from
\eqref{orb-related-lattices}.  Applying the inverse shears
shows that the resulting transformation is a bijection of
the finite set of pairs grouped under $\Psi$.

Only one adjacent upper entry occurs in the shears, the
$(k,k+1)$ entry.  Since $\beta_{k+1}=\beta'_{k+1}=1$, the
change in the difference of the upper phases is
\[
 \xi_k-\xi'_k=-r_0d.
\]
Thus this permutation multiplies every weight in the finite
sum for $\Psi$ by the same number $\psi(-r_0d)\ne1$.
That sum must be zero.  Summing over the finitely many lifts
and over the outer cosets proves the proposition.

Finally, using canonical representatives $p_B$ has not
changed the original counts.  A right integral upper
unipotent change replaces the inverse-coordinate functionals
by an integral upper unipotent combination.  It preserves
the last functional and every condition $T_j$, in both
source and target.  Its phase is trivial.  Hence all the
partitions and changes of representatives used above are
legitimate in \eqref{orb-actual-kernel}.
\end{proof}

\begin{corollary}\label{orb-support-proved}
For the actual kernel \eqref{orb-actual-kernel}, if $s\ge1$,
the previous block configurations $h,h'$ have equal total
size, and
$\mathcal K(h[s,t],h'[s,t'])\ne0$, then
\[
 t^{-1}-t'^{-1}\in R.
\]
\end{corollary}

\begin{proof}
Put $r=\dim Z(h)=\dim Z(h')$ and $n=r+s$.
Marked isomorphisms already satisfy $T_{n-1}$.
Repeated use of Proposition~\ref{orb-triangular-average}
shows that the same kernel is obtained by counting only
isomorphisms satisfying $T_r$.
If its value is nonzero, at least one such isomorphism occurs.
For $Z=xZ(h[s,t])y$,
\[
 Z^{-1}e_{n-1}=t^{-1}y^{-1}e_r.
\]
Its $r$th coordinate is $t^{-1}$ and all coordinates greater
than $r$ are zero.  The corresponding coordinates for $Z'$
are $t'^{-1}$ and zeros.  Evaluate the relation
\eqref{orb-triangular} with $i=r$ on the marked vector,
which the isomorphism preserves.  It gives
$t'^{-1}\equiv t^{-1}\pmod R$, as asserted.
When $s=1$ no averaging step is needed: $r=n-1$ already.
\end{proof}

\paragraph{Exact finite controls.}
An accompanying integer-arithmetic program checks the averaging
identity for a single block of size two over $\mathbb Q_p$,
with $p=3,5$, $t=p$, and $t'=vp$ for $v=1,2,p+1$.
For a block with scalar $pv$, writing the upper unipotent
parameter as $b/p$ with $0\le b<p$ gives the integral orbit matrix
\[
 Z_b=\begin{pmatrix}0&pv\\pv&vb\end{pmatrix}.
\]
All marked-module isomorphisms are enumerated directly in
$(\mathbb Z/p^2\mathbb Z)^2/Z_b(\mathbb Z/p^2\mathbb Z)^2$,
using the possible images of its two generators.  The program
independently tests the additional condition $T_0$ and reduces
the phase sums modulo $1+X+\cdots+X^{p-1}$.
For $v=2$, the unrestricted phase polynomial is
$p(1+X+\cdots+X^{p-1})$, whereas the restricted polynomial is
zero: cancellation is essential.  For $v=1,p+1$, both values
are $p^2$.  These six exact controls agree with the support
corollary.  Four additional exact congruence checks, for
$N=4,5,6,7$, verify the representative change in the preceding
$\mathbb Q_3$ counterexample.  The general assertions rest on
the displayed proofs, not on these finite computations.

\paragraph{Scope of the reconstructed truncated orbital identity.}
For the actual finite marked-cokernel kernel, its Gram formula
proves Hermitian symmetry, Proposition~\ref{orb-local-constancy-proved}
proves the required local constancy uniformly against all comparison
configurations, and Corollary~\ref{orb-support-proved} proves the
inverse-parameter support.  The following finite-type argument
reduces comparison-configuration uniformity to finitely many
fixed marked modules.  Section~\ref{hr-section} constructs the
actual matrix test functions, and
Section~\ref{sec:theta-scalar-germ} proves their algebraic
reduction to the oscillatory integral.  Together with the previously
proved stationary-phase formula, these give the actual kernel germ
in Theorem~\ref{ga-actual-germ}.  Thus all hypotheses of
Theorem~\ref{orb-cut-identity} are established for this kernel.
The averaged diagonal estimate needed for the positive-majorant
induction is a separate finite character calculation, proved in
Section~\ref{tp-section}.  The quantitative global theta energy
estimates remain separate statements; they do not follow merely
from these local identities.


\subsection{Reduction of the germ formula to finitely many fixed marked modules}
\label{gu-section}

The comparison configuration in the remaining orbital germ formula is
arbitrary. Its uniformity can nevertheless be reduced to finitely many
fixed targets. The reason is that the determinant stays in a compact
subset of $K^\times$. This observation does not prove the scalar germ
formula, but removes a separate uniformity requirement from that task.
We retain the local field product $K$, compact open subring $R$, and
marked cokernels used in Section~\ref{sec:theta-orbital-identity}.

\begin{lemma}\label{gu-bounded-types}
Let $n\ge0$ and let $D\subset K^\times$ be compact. There is an
integer $N\ge0$, depending only on $R$ and $D$, such that every
$Z\in M_n(R)\cap\mathrm{GL}_n(K)$ with $\det Z\in D$ satisfies
\[
             p^N R^n\subset ZR^n\subset R^n.
\]
Consequently the finite modules $C_Z=R^n/ZR^n$, with any specified
vector $e\in C_Z$ and any $R$-linear functional
$\lambda:C_Z\to K/R$, belong to finitely many isomorphism classes.
Writing $a=\#(R/p^NR)$ and $B=a^n$, their number is at most
\[
                         2^B B a^B.
\]
These assertions hold for nonmaximal compact open orders as well.
\end{lemma}

\begin{proof}
The inverse set $D^{-1}$ is compact. Since $R$ is a neighborhood
of zero, multiplication by a sufficiently large power of $p$ sends
all of $D^{-1}$ into $R$; choose such a power $p^N$.
The adjugate of an integral matrix has entries in $R$. Therefore
\[
 p^N Z^{-1}=\frac{p^N}{\det Z}\operatorname{adj}(Z)
                                      \in M_n(R),
\]
which proves the first inclusion. The second is integrality of $Z$.
The case $n=0$ has the zero module and also satisfies the conclusion.

The group $R/p^NR$ is finite: it is compact and discrete, because
$p^NR$ is an open subgroup of $R$. Put $A=R/p^NR$.
The module $C_Z$ is a quotient of the finite module $A^n$ and has
cardinality at most $B$. There are at most $2^B$ choices for its
defining submodule, at most $B$ choices for its marked vector,
and at most $a^B$ choices for its functional. For the last claim,
$p^NC_Z=0$ implies that the functional takes values in
\[
          (K/R)[p^N]=p^{-N}R/R,
\]
which has cardinality $a$ by multiplication by $p^N$.
Counting all functions into this finite set only increases the
number of $R$-linear ones. These finite presentations cover all
the marked modules under consideration, so also bound their
number of isomorphism classes. No principal-ideal property of
$R$ is used.
\end{proof}

For a finite $R$-module with specified vector and functional,
write $T=(C,e,\lambda)$ and define its orbital coefficient by
\begin{equation}\label{gu-orbital-coefficient}
 \mathcal B(w;T)=
       \sum_{q\in\mathcal I_w}a_q\,
                     \#\operatorname{Iso}(C_{Z_q},T).
\end{equation}
Isomorphisms preserve both specifications. This is a finite sum by
Lemma~\ref{orb-finite-orbits}. It depends only on the isomorphism
class of $T$. Directly from the definition of the kernel,
\begin{equation}\label{gu-kernel-coefficients}
 \mathcal K(w,w')=
       \sum_{q'\in\mathcal I_{w'}}\overline{a_{q'}}\,
                        \mathcal B(w;C_{Z_{q'}}).
\end{equation}
The notation on the right includes the specified vector and
functional of each cokernel.

\begin{proposition}\label{gu-uniform-germ}
Fix $m\ge1$, put $s=m+1$, and fix a finite collection
$\mathcal F$ of configurations whose total size is $r$.
Put $n=r+s$ and fix a compact set $Q\subset K^\times$.
For a critical root $z^s=(-t)^mc$, write
\[
 b(t,c;z)=|t|^{-m^2/2}|s|^{-1/2}
              \gamma_m(zt)|z|^{m^2/2}\psi(sz/t).
\]
Suppose that for each $h\in\mathcal F$ and each fixed
$Z'\in M_n(R)\cap\mathrm{GL}_n(K)$ there is a neighborhood
$V_{h,Z'}$ of zero such that, whenever $t,c\in K^\times$,
$t\in V_{h,Z'}$ and $ct^m\in Q$, one has
\begin{equation}\label{gu-pointwise-germ}
 \mathcal B(W_h(t,c);C_{Z'})
       =\sum_{z^s=(-t)^mc}
                 b(t,c;z)\mathcal B(h[s,z];C_{Z'}).
\end{equation}
Then there is one neighborhood $V$ of zero such that, for every
$h\in\mathcal F$, $t,c\in K^\times$ with $t\in V$ and
$ct^m\in Q$, and every comparison configuration $w'$, one has
\begin{equation}\label{gu-uniform-germ-conclusion}
 \mathcal K(W_h(t,c),w')
       =\sum_{z^s=(-t)^mc}
                    b(t,c;z)\mathcal K(h[s,z],w').
\end{equation}
The total size and parameters of $w'$ are unrestricted. In fact,
only finitely many fixed matrices $Z'$ are needed in the
hypothesis.
\end{proposition}

\begin{proof}
Let $H=Z(h)$. The determinant of a reversing matrix of size $j$
is $(-1)^{j(j-1)/2}$. Hence
\[
 \det Z(W_h(t,c))
      =\det H\,(-1)^{m(m-1)/2}t^mc.
\]
For every critical root,
\begin{align*}
 \det Z(h[s,z])
  &=\det H\,(-1)^{m(m+1)/2}z^{m+1}\\
  &=\det H\,(-1)^{m(m+1)/2+m}t^mc
    =\det Z(W_h(t,c)).
\end{align*}
The last two exponents differ by $2m$. Multiplication by the
lower and upper unipotent matrices in an orbital term does not
change the determinant.

Thus the determinants of all integral matrices occurring on
either side, for every allowed parameter and every
$h\in\mathcal F$, lie in one compact subset $D\subset K^\times$.
Lemma~\ref{gu-bounded-types} gives a finite list of possible marked
cokernel types. Keep only those which actually occur on one of
the two sides for some allowed parameter, and choose a realizing
integral matrix $Z_T\in M_n(R)\cap\mathrm{GL}_n(K)$ for each type
$T$. This list and its representatives depend on the already
fixed $\mathcal F,m,Q$, not on any later $t,c$ or $w'$.
If the list is empty, all relevant orbital coefficients vanish
and the result is immediate.

Intersect the finitely many neighborhoods $V_{h,Z_T}$.
For $t$ in their intersection, \eqref{gu-pointwise-germ} holds
for every possible target type $T$. It also holds with any target
of a different type: both sides are zero, because an isomorphism
can exist only between equal marked types.

Apply these identities to the finitely many marked modules
$C_{Z_{q'}}$ occurring in an arbitrary comparison configuration
$w'$. Multiply by $\overline{a_{q'}}$ and sum. Formula
\eqref{gu-kernel-coefficients} and interchange of the finite sums
give \eqref{gu-uniform-germ-conclusion}. The chosen neighborhood
is independent of $w'$, as required.
\end{proof}

The proposition supplies precisely the comparison-configuration
uniformity in \eqref{orb-germ-input}. The scalar identity
\eqref{gu-pointwise-germ} is supplied by the initial-block
reduction in Section~\ref{hr-section} and the scalar germ in
Section~\ref{sec:theta-scalar-germ}. The present proposition
shows that, for fixed $m$, $\mathcal F$ and $Q$,
one need not establish bounds simultaneously for infinitely many
unrelated target modules. The finite list is determined before
the small parameter is chosen, so no exchange of an infinite
intersection with a neighborhood occurs.


